\documentclass[12pt]{article}
%\AddToHook{cmd/section/before}{\clearpage} % To start each section in a new page
\usepackage[a4paper,
            bindingoffset=0in,
            left=1in,
            right=1in,
            top=1in,
            bottom=1in,
            footskip=.5in]{geometry}
\usepackage{epsfig}
\usepackage{amsmath,graphicx,psfrag,pstricks} % Using pstricks instead of pstcol (depecrated)
\usepackage{amssymb} % Math symbols
\usepackage[font=small,labelfont=bf, position=top]{caption} % Required for specifying captions to tables and figures
\usepackage{subcaption} % Combine figures  
\usepackage[section]{placeins} % Force figures to stay in their section
\usepackage{minted} % For inserting formated Python code
\usepackage{listings} % For inserting code from files
\def\n{\noindent}
\def\u{\underline}
\def\hs{\hspace}
\newcommand{\thrfor}{.^{\displaystyle .} .}
\newcommand{\bvec}[1]{{\bf #1}}
\usepackage{hyperref} % To typeset URLs
\usepackage[sorting=none,style=ieee]{biblatex} %Imports biblatex package. Sorting none, so that number in order of appearance
\addbibresource{GG2.bib} %Import the bibliography file
\usepackage{listings}
\usepackage{xcolor} % Optional for syntax highlighting
\usepackage[toc,page]{appendix}
\usepackage{booktabs}

\lstset{ 
  language=C++,                      % Language of the code
  basicstyle=\ttfamily\small,         % Font and size of the code
  keywordstyle=\color{blue}\bfseries, % Style for keywords
  commentstyle=\color{green},         % Style for comments
  stringstyle=\color{red},            % Style for strings
  frame=single,                       % Frame around the code
  numbers=left,                       % Line numbers
  numberstyle=\tiny\color{gray},      % Style for line numbers
  stepnumber=1,                       % Number every line
  numbersep=5pt,                      % Space from code to numbers
  breaklines=true                     % Break long lines
}


\begin{document}

\noindent
\rule{15.7cm}{0.5mm}


\begin{center}
{\Large \bf ENGINEERING TRIPOS PART IIB \\
\bigskip}
\vspace{2cm}
{\LARGE \bf Modelling the Olfactory Bulb\\
\bigskip}
{\Large \bf  Technical Milestone Report \\\hfill \\}
\bigskip
\textsc{\Large Hana Iza Kim} \\\hfill\\
{\large Pembroke College, January 2026}\\
\hfill\\
\end{center}
\rule{15.7cm}{0.5mm}
\vfill
\bigskip
\vfill
\thispagestyle{empty} % Skip numbering the first page
\pagebreak
\setcounter{page}{1}


\newpage

%%%%%%%%%%%%%%%%%%%%%%%%%%%%%%%%%%%%%%%%%%%%%%%%%%%%%%%%%%
\section{Introduction}
%%%%%%%%%%%%%%%%%%%%%%%%%%%%%%%%%%%%%%%%%%%%%%%%%%%%%%%%%%

Homeostatic plasticity is a fundamental mechanism by which neural circuits maintain stable function despite perturbations. This process allows neurons to adapt their properties while preserving overall circuit output, enabling robust information processing in the face of environmental changes. Understanding the mechanisms underlying homeostatic plasticity is crucial for both basic neuroscience and potential therapeutic applications.

The olfactory bulb (OB) in mice provides an ideal model system for investigating homeostatic plasticity. This circuit exhibits remarkable resilience: when sensory input is temporarily deprived, structural changes occur (such as axonal shrinkage in dopaminergic neurons), yet functional output remains unchanged upon restoration of input. This phenomenon suggests that compensatory changes must occur elsewhere in the circuit to maintain output stability.

The central research question addressed in this work is: \textit{What parameter changes in the olfactory bulb circuit allow structural plasticity while preserving functional output?} To answer this, we develop a mathematical model of the OB circuit, identify parameter configurations that produce target outputs, analyze the stability of these configurations, and characterize the space of parameter changes that preserve output—thereby providing insight into the mechanisms of homeostatic plasticity.

The objectives of this technical milestone are: (1) to develop a simplified but biologically motivated mathematical model of the OB circuit, (2) to find parameter configurations that satisfy known input-output relationships using gradient-based optimization, (3) to analyze the stability of these solutions, and (4) to identify directions in parameter space along which changes preserve output, corresponding to potential mechanisms of homeostatic plasticity.

\section{Homeostatic plasticity}
Homeostatic plasticity refers to the ability of neural circuits to maintain stable function despite perturbations through adaptive changes in neuronal properties. The olfactory bulb in mice provides a convenient model system to investigate this phenomenon.

The key experimental observation is as follows: when an odour is presented to a mouse, the output signal to the brain is recorded. The mouse's nose is one-sided, so blocking the right nostril only affects the right half of the brain, and similarly for the left side. When sensory deprivation is induced by placing a nose plug, it was observed that dopaminergic axonic sensory neurons shrink in axonic length \cite{galliano2021}. This structural change is analogous to muscle atrophy when a limb is immobilized in a cast. However, when the nose plug is removed, the signal output from the OB returns to the same level as before sensory deprivation. This remarkable observation suggests that compensatory changes must occur elsewhere in the neuron circuit such that the output of the OB remains unchanged despite the structural alterations. 

\begin{figure}[H]
    \centering
    \includegraphics[width=0.8\textwidth]{Figures/experiment.png}
    \caption{Question: what neurological changes can occur in the olfactory bulb that allow plasticity while preserving functionality?}
    \label{fig:experiment}
\end{figure}

\section{Olfactory bulb circuit}
The circuitry of the olfactory bulb forms the foundation of our model. Each odour is a chemical composition that binds to a specific combination of receptors. Interestingly, each receptor type projects to the same glomeruli, which are regions where neurotransmitters are released. Within the glomeruli, many interneurons are at play. We group the interneurons into axonic and anaxonic interneurons, and external tufted cells. 

\begin{itemize}
  \item \textbf{Blue neurons}: Inhibitory neurons that act as self-inhibitors, preventing over-excitation
  \item \textbf{Green neurons}: Inhibitory neurons that receive input from one glomerulus and inhibit many neighbouring glomeruli- their axons are very long and clearly visible.
\end{itemize}

External tufted cells (ETCs) play a crucial role in the circuit architecture. In a mouse's brain, there are two olfactory bulbs, each corresponding to the right and left side of the brain. Within each olfactory bulb, there are two glomeruli of each type of receptor cell. The external tufted cells connect the two sister cells of the same type of glomeruli and send a signal to the corresponding mitral cell. (For the purpose of the model, we merge the sister cells together, since we are simplifying the real system.)

The mitral cell integrates signals from the glomeruli and sends an output signal to the cortex. 

The final layer is the granule layer, where granule cells receive input from both the cortex and the mitral cell, and provide inhibitory output back to the mitral cell. This acts as lateral inhibition similar to the green cells, but is more connected to learning processes due to its cortical input.

\begin{figure}[H]
    \centering
    \includegraphics[width=0.8\textwidth]{Figures/OB_1.png}
    \caption{Pink/orange/brown= Olfactory Sensory Neuron (OSN) ; green =
            Inhibitory axonic interneurons ; blue = Inhibitory anaxonic interneuron ; red = External Tufted Cells (ETCs) ; grey = Mitral cells ; black = Granule Cells}
    \label{fig:OB_1}
\end{figure}




\section{Mathematical interpretation of the model}
Let us consider a simplified model where input passes through a nonlinear activation function. We choose the Naka-Rushton function, which is defined as:
\[
f(x)=\frac{x^n}{x^n +k^n}
\]

The input to the next neuron is a fraction of the output of the the previous neuron, this can be represented by the weight. Note, these signals represent firing rate and the weights represent synaptic strength.

By letting $x'=x/k$, the scalar $k$ can be absorbed into the weight values, leading to the simplified Naka-Rushton expression:
\[
f(x')=\frac{x'^n}{x'^n + 1}
\]
With $n=4$, this form saturates at 1 for large inputs. For our model, we aim to keep neuron output firing rates around 0.5 to operate in the nonlinear region of the function, where the response is most sensitive to input changes. 

\begin{figure}[H]
    \centering
    \includegraphics[width=0.6\textwidth]{Figures/simple_circuit.png}
    \caption{Simply small model, x, y , z represent firing rate and weight represents synaptic strength. }
    \label{fig:simple_circuit}
\end{figure}

\begin{figure}[H]
    \centering
    \includegraphics[width=0.6\textwidth]{Figures/naka_rushton.png}
    \caption{Naka-Rushton }
    \label{fig:naka_rushton}
\end{figure}






\section{Connectivity - simplifying the circuit}
Having understood the key players, we examine the connectivity and determine which neurons interact with each other. The simplified model makes the following assumptions:

\begin{enumerate}
    \item \textbf{Cell grouping}: We group together "types" of cells and represent each type by a single representative cell.
    \item \textbf{Two glomeruli}: We assume only two glomeruli (to enable demonstration of cross-glomerular inhibition by green cells).
    \item \textbf{Self-inhibition}: Blue neurons inhibit themselves.
    \item \textbf{No temporal effects}: We neglect temporal dynamics and focus on steady-state behavior.
    \item \textbf{Steady-state signals}: We consider only steady-state firing rates. While there are oscillations in firing rate due to sniffing and high-frequency oscillations in the real system, this assumption is valid because we are comparing end states rather than transients. The steady-state approximation captures the long-term behavior relevant to homeostatic plasticity.
\end{enumerate}

\begin{figure}[H]
    \centering
    \includegraphics[width=0.6\textwidth]{Figures/OB_circuit.png}
    \caption{OB circuit simplified, signals represent firing rate and weight represents synaptic strength. }
    \label{fig:OB_circuit}
\end{figure}

The circuit dynamics are represented by the following equations:

\begin{align}
b &= \operatorname{nr}\!\left(w_{\text{osn},b}\,\alpha x\right) \label{eq:b} \\[6pt]
r &= \operatorname{nr}\!\left(
    w_{\text{osn},r}\,\alpha x
    - w_{b,r}\, b
    - \mathbf{W}_{g,r}\, g
\right) \label{eq:r} \\[6pt]
g &= \operatorname{nr}\!\left(
    w_{\text{osn},g}\,\alpha x
    - w_{b,g}\, b
    + \mathbf{W}_{r,g}\, r
\right) \label{eq:g} \\[6pt]
y &= \operatorname{nr}\!\left(
    w_{\text{osn},y}\,\alpha x
    - w_{b,y}\, b
    - \mathbf{W}_{g,m}\, g
    + w_{r,y}\, r
\right) \label{eq:y} \\[10pt]
g_r &= \operatorname{nr}\!\left(
    w_{\text{cor}}\, \text{cortex}
    + w_y\, y
\right) \label{eq:gr} \\[6pt]
y &= y - w_{g_r,y}\, g_r \label{eq:y_feedback}
\end{align}

Since we are considering steady-state values, there is a feedback loop from OSN to blue neurons with an inhibitory signal back to OSN. Assuming this results in a steady-state output, we represent this steady-state value as a fraction of the original signal $x$, denoted as $\alpha x$, where $\alpha \in [0,1]$ captures the effect of the feedback loop.

Since g only influences neurons of the other receptor type, matrices are used. 

Notice that the last two equations (Equations~\eqref{eq:gr} and \eqref{eq:y_feedback}) are only valid if we include granule cells in our model. Although granule cells play an important role in tuning the output $y$ of the mitral cell, they depend on the cortical input signal, which is not well characterized experimentally. Granule cells are also more associated with learning processes than with the homeostatic mechanisms of interest here. Since granule cells simply result in another fixed-point iteration problem and we lack sufficient data to constrain their parameters, we exclude them from the following analysis and parameter optimization. 

\section{Fixed iteration problem}
From the above equations, it can be seen that solving these equations results in a fixed-point iteration problem. Specifically, Equation~\eqref{eq:g} can be rewritten by substituting Equations~\eqref{eq:b} and \eqref{eq:r}:

\[
g = \operatorname{nr}\!\left(
    w_{\text{osn},g}\,\alpha x
    - w_{b,g}\, \operatorname{nr}\!\left(w_{\text{osn},b}\,\alpha x\right)
    + \mathbf{W}_{r,g}\, (\operatorname{nr}\!\left(
    w_{\text{osn},r}\,\alpha x
    - w_{b,r}\, \operatorname{nr}\!\left(w_{\text{osn},b}\,\alpha x\right)
    - \mathbf{W}_{g,r}\, g)
\right)
\right) \\[6pt]
\]
and g is a function of itself. 

To solve this fixed-point equation, we employ Anderson Acceleration, a method that accelerates the convergence of standard fixed-point iteration. Instead of the naive approach where $g_{k+1} = f(g_k)$ iteratively updates the guess using only the previous iterate, Anderson Acceleration maintains a history of the $m$ most recent iterates and their corresponding residuals $r_k = f(g_k) - g_k$. At each iteration, the algorithm computes the next guess by solving a least-squares problem to find optimal mixing coefficients $\alpha_i$ that minimize the residual when forming a linear combination: $g_{k+1} = \sum_{i=0}^{m} \alpha_i f(g_{k-i})$. This extrapolation leverages information from multiple previous iterations to predict a better next guess, effectively approximating the inverse of the Jacobian and accelerating convergence, particularly when the fixed-point mapping is contractive but converges slowly. In this implementation, we configured Anderson Acceleration with a history size of 5, tolerance of $10^{-6}$, and a maximum of 50 iterations, which efficiently converges to the fixed-point solution for the green neuron firing rates.



\section{Parameter setting}
Given a model of the neural circuit structure, and assuming we can measure the input odour signal and the output signal of the OB, we are interested in finding possible weight configurations in our model that satisfy these known input-output relationships.

The strategy to achieve this is as follows: we sample initial points in parameter space and then use the Adam optimizer to find the nearest parameter configuration that produces the target output.

The loss function to be minimized is defined as:
\begin{equation}
L(\theta) = \sum_{i=1}^{m} \left[(y_i - y_{\text{target},i})^2 + (g_i - g_{\text{target},i})^2 + (b_i - b_{\text{target},i})^2 + (r_i - r_{\text{target},i})^2\right]
\label{eq:loss}
\end{equation}
where $m$ is the number of mitral neurons (channels), and $y_i$, $g_i$, $b_i$, $r_i$ are the firing rates of the mitral, green, blue, and red neurons respectively, computed from the model with parameters $\theta$.

The Adam optimizer uses exponential moving averages of past gradients (first moment with $\beta_1=0.9$) and squared gradients (second moment with $\beta_2=0.999$) to adaptively update learning rates. Our implementation uses JAX/optax with JIT compilation for efficient gradient computation via automatic differentiation. The optimization loop enforces constraints by: (1) clipping parameters to be non-negative, and (2) setting diagonal elements of weight matrices to zero (since neurons do not self-connect in our model). The learning rate decays if no improvement is observed for a specified number of consecutive evaluations.

During optimization, we constrain output firing rates to be positive and approximately 0.5. This ensures that neurons operate in the nonlinear region of the Naka-Rushton function, where the response is most sensitive to input changes. We then analyze the stability of each optimized parameter set. 

\begin{figure}[H]
    \centering
    \begin{subfigure}[b]{0.48\textwidth}
        \centering
        \includegraphics[width=\textwidth]{Figures/adam_opt_loss_and_output_y.png}
        \caption{}
        \label{fig:adam_opt}
    \end{subfigure}
    \hfill
    \begin{subfigure}[b]{0.48\textwidth}
        \centering
        \includegraphics[width=\textwidth]{Figures/rgby_outputs_after_adam_opt_of_param.png}
        \caption{}
        \label{fig:rgby_outputs}
    \end{subfigure}
    \caption{Results of Adam optimization: (a) training loss trajectory and output $y$ during optimization; (b) resulting R, G, B, and Y outputs after parameter optimization via Adam opt.}
    \label{fig:adam_opt_combined}
\end{figure}


\section{Stability}

To analyze the stability of the network around a candidate fixed point, we consider the state vector $s(t)$ representing the outputs of the neurons across all channels. Let $s^\ast$ denote the steady-state solution for a given set of parameters $\theta$, i.e., the fixed point where 
\[
F(s^\ast, \theta, x, m) = s^\ast.
\]

We introduce a small perturbation $e(t)$ around the fixed point:
\[
s(t) = s^\ast + e(t),
\]

Linearizing the dynamics around $s^\ast$ gives
\[
\tau \frac{d e}{dt} = - e + J \, e, \quad \text{where} \quad
J = \left. \frac{\partial F}{\partial s} \right|_{s = s^\ast}.
\]
Rewriting this, we obtain a standard continuous-time linear system:
\[
\frac{d e}{dt} = \frac{1}{\tau} (J - I) \, e = A \, e.
\]
The fixed point is locally asymptotically stable if all eigenvalues of $A$ have negative real parts, ensuring that the perturbation $e(t)$ decays over time.

For the optimized parameter sets, we compute the Jacobian matrix $J$ at the fixed point using automatic differentiation, construct the stability matrix $A = \frac{1}{\tau}(J - I)$, and check whether all eigenvalues have negative real parts. This analysis reveals which parameter configurations yield stable network dynamics, a necessary condition for biologically plausible solutions.

In future analysis, we plan to incorporate more biological insights, such as the ratio of excitatory to inhibitory input to each neuron, or the relative strengths between different neuron types. 

\section{Parameter changes preserving output ($\Delta \theta$)}

Given a set of biologically realistic parameters $\theta$, we seek to understand how these parameters could change such that the output of the system remains unchanged. This directly addresses the initial problem: during sensory deprivation, structural changes occur (e.g., axonal shrinkage), yet functional output is preserved, suggesting compensatory parameter changes elsewhere in the circuit.

To find parameter changes $\Delta \theta$ that preserve output, we require $y(\theta + \Delta \theta) = y(\theta)$. Taking a Taylor expansion:
\[
y(\theta + \Delta \theta) \approx y(\theta) + g^T \Delta \theta + \frac{1}{2} \Delta \theta^T H \, \Delta \theta + \mathcal{O}(\|\Delta \theta\|^3)
\]
where $g = \nabla_\theta y(\theta)$ is the gradient and $H = \nabla_\theta^2 y(\theta)$ is the Hessian matrix. For the output to remain unchanged, we require:
\begin{equation}
y(\theta + \Delta \theta) - y(\theta) = g^T \Delta \theta + \frac{1}{2} \Delta \theta^T H \, \Delta \theta \approx 0.
\label{eq:delta_y}
\end{equation}

This condition can be satisfied in several ways: (1) $\Delta \theta$ is perpendicular to the gradient ($g^T \Delta \theta = 0$) and lies in the null space of the Hessian, (2) $\Delta \theta$ is an eigenvector of $H$ corresponding to a zero eigenvalue, or (3) the first- and second-order terms cancel each other out (though this typically yields isolated solutions rather than a continuous manifold).

To characterize the space of output-preserving parameter changes, we compute the Hessian matrix $H$ at an optimized parameter set $\theta^*$ using automatic differentiation. The Hessian is computed as $H = \nabla_\theta^2 L(\theta^*)$, where $L(\theta)$ is the loss function defined in Equation~\eqref{eq:loss}. We then perform an eigendecomposition: $H = U \Lambda U^T$, where $\Lambda$ is a diagonal matrix of eigenvalues and $U$ contains the corresponding eigenvectors.

Analysis of the eigenvalue spectrum reveals that there are several eigenvalues that are very small (near zero). This is promising, as it suggests the existence of a manifold of parameter changes along which the output remains approximately constant. The eigenvectors corresponding to these near-zero eigenvalues define directions in parameter space where changes have minimal impact on the output.

\begin{figure}[H]
    \centering
    \includegraphics[width=0.8\textwidth]{Figures/hessian_eigenvalue_spectrum.png}
    \caption{Hessian eigenvalue spectrum showing several near-zero eigenvalues, indicating directions in parameter space where changes preserve output.}
    \label{fig:hessian_spectrum}
\end{figure}

\begin{figure}[H]
    \centering
    \includegraphics[width=0.8\textwidth]{Figures/hessian_eigenvectors_small.png}
    \caption{Eigenvectors corresponding to the four smallest eigenvalues of the Hessian matrix. These eigenvectors encode the structure of compensatory parameter changes that preserve output.}
    \label{fig:hessian_eigenvectors}
\end{figure}

Specifically, if $\lambda_i \approx 0$ is a near-zero eigenvalue with corresponding eigenvector $u_i$, then moving along $\Delta \theta = \epsilon u_i$ (for small $\epsilon$) will approximately preserve the output, since $\Delta \theta^T H \Delta \theta = \epsilon^2 \lambda_i \approx 0$. The participation ratio, defined as $(\sum_i \lambda_i)^2 / \sum_i \lambda_i^2$, quantifies the effective dimensionality of the parameter space relevant to output changes.

The eigenvectors corresponding to the smallest eigenvalues encode the structure of compensatory changes. For example, if a particular eigenvector shows increased weights in one pathway and decreased weights in another, this suggests a compensatory mechanism: increased activity in one component is balanced by decreased activity in another, maintaining overall output. The challenge is to interpret these eigenvectors in biological terms, relating them to known mechanisms of homeostatic plasticity such as synaptic scaling, changes in connection strength, or alterations in interneuron activity.

\subsection{Biological Interpretation of Compensatory Parameter Changes}

To understand how parameter changes preserve output in biological terms, we analyze the eigenvectors corresponding to the smallest eigenvalues of the Hessian matrix. These eigenvectors reveal the structure of compensatory mechanisms: when certain parameters increase, others must decrease to maintain constant output.

Consider the eigenvector $u_i$ corresponding to a near-zero eigenvalue $\lambda_i \approx 0$. This eigenvector has components corresponding to each parameter in $\theta$. By examining which components are positive (indicating parameter increases) and which are negative (indicating parameter decreases), we can identify compensatory relationships.

For example, if an eigenvector shows:
\begin{itemize}
    \item Positive components for weights $w_{OSN \to r}$ (OSN to red neuron connections) and $W_{r \to g}$ (red to green connections)
    \item Negative components for weights $w_{OSN \to g}$ (OSN to green neuron connections) and $W_{g \to r}$ (green to red connections)
\end{itemize}

This pattern suggests a compensatory mechanism: if OSN input to red neurons increases, and red-to-green excitation increases, then to maintain constant output, OSN input to green neurons must decrease, and green-to-red inhibition must decrease. This creates a balance where increased activity in the red pathway is compensated by decreased activity in the green pathway.

Such compensatory relationships are consistent with known mechanisms of homeostatic plasticity. During sensory deprivation, when one pathway experiences reduced input, other pathways may strengthen to compensate. The eigenvectors of the Hessian matrix reveal which specific parameter combinations allow this compensation while preserving output.

By analyzing the eigenvectors corresponding to the two smallest eigenvalues, we can identify the primary compensatory mechanisms in the circuit. These eigenvectors define the directions in parameter space along which the system can adapt during sensory deprivation while maintaining functional output.

\section{Results}

\subsection{Parameter Optimization Results}

Using the Adam optimizer with the loss function defined in Equation~\eqref{eq:loss}, we successfully found multiple parameter configurations that produce target firing rates of approximately 0.5 for all neuron types. Starting from 10 randomly sampled initial parameter sets, the optimization converged to solutions with final loss values on the order of $10^{-6}$ or smaller. The optimized outputs closely match the target values, with firing rates typically within $10^{-4}$ of the target value of 0.5.

The optimization results demonstrate that there exist multiple distinct parameter configurations that satisfy the same input-output relationship. This multiplicity suggests that the system has redundancy, which is a necessary condition for homeostatic plasticity: if multiple parameter sets can produce the same output, then the system can adapt its parameters while preserving function.

\subsection{Stability Analysis Results}

Stability analysis was performed on all optimized parameter sets. For each solution, we computed the Jacobian matrix at the fixed point and determined the eigenvalues of the stability matrix $A = \frac{1}{\tau}(J - I)$. The results show that a significant fraction of the optimized parameter sets yield stable fixed points, with all eigenvalues having negative real parts. This indicates that these solutions are not only functionally correct (producing the target output) but also dynamically stable, making them biologically plausible.

The eigenvalue distributions reveal the timescales of perturbation decay, with more negative real parts corresponding to faster recovery to the fixed point. This information could be used to constrain parameter values further, ensuring that the network dynamics match observed biological timescales.

\begin{figure}[H]
    \centering
    \includegraphics[width=0.8\textwidth]{Figures/stability_analysis.png}
    \caption{Stability analysis results: (a) real parts of eigenvalues for all samples; (b) maximum real part indicating stability; (c) eigenvalue spectrum in complex plane; (d) summary statistics.}
    \label{fig:stability_analysis}
\end{figure}

\subsection{Hessian Analysis Results}

The Hessian matrix analysis reveals the structure of the parameter space around optimized solutions. The eigenvalue spectrum of the Hessian shows a clear separation: most eigenvalues are significantly larger than zero, indicating directions in parameter space where changes strongly affect the output, while a smaller number of eigenvalues are near zero, indicating directions where changes have minimal impact on output.

The participation ratio, computed from the Hessian eigenvalues, quantifies the effective dimensionality of the output-relevant parameter space. A low participation ratio (relative to the total number of parameters) indicates that the output depends on only a subset of parameter combinations, leaving room for compensatory changes in other directions.

\begin{figure}[H]
    \centering
    \includegraphics[width=0.8\textwidth]{Figures/participation_ratio.png}
    \caption{Participation ratio analysis: (a) participation ratio across multiple parameter sets, showing effective dimensionality; (b) number of near-zero eigenvalues indicating output-preserving directions.}
    \label{fig:participation_ratio}
\end{figure}

The eigenvectors corresponding to the smallest eigenvalues define the directions of output-preserving parameter changes. Analysis of these eigenvectors reveals patterns of coordinated changes: for example, increases in certain synaptic weights are paired with decreases in others, suggesting compensatory mechanisms that maintain overall circuit function.

\section{Conclusion}

This technical milestone report presents a mathematical framework for understanding homeostatic plasticity in the olfactory bulb circuit. We developed a simplified but biologically motivated model of the OB circuit, incorporating key neuron types (blue, green, red interneurons, external tufted cells, and mitral cells) and their connectivity patterns. The model uses the Naka-Rushton activation function to capture nonlinear neural responses and operates in steady-state to focus on long-term adaptive changes.

Using gradient-based optimization (Adam optimizer), we identified multiple parameter configurations that produce target firing rates, demonstrating the existence of parameter redundancy—a necessary condition for homeostatic plasticity. Stability analysis confirmed that many of these solutions are dynamically stable, making them biologically plausible.

The key finding is that the Hessian matrix of the loss function has a significant number of near-zero eigenvalues, indicating the existence of a manifold of parameter changes that preserve output. This provides a mathematical foundation for understanding how the circuit can adapt its parameters (e.g., during sensory deprivation) while maintaining functional output.

\subsection{Limitations and Future Work}

Several limitations should be acknowledged. First, the model is highly simplified, grouping neuron types and neglecting temporal dynamics. Second, we do not include granule cells in the main analysis due to insufficient data on cortical input. Third, the steady-state assumption ignores important oscillatory dynamics present in the real system. Fourth, the biological interpretation of the Hessian eigenvectors requires further validation against experimental data.

Future work should: (1) incorporate more detailed biological constraints (e.g., excitatory-inhibitory balance ratios), (2) extend the model to include temporal dynamics and oscillations, (3) validate predicted parameter changes against experimental measurements during sensory deprivation, (4) analyze the biological meaning of the output-preserving directions in parameter space, and (5) investigate how these mechanisms relate to known forms of homeostatic plasticity such as synaptic scaling and structural plasticity.

The framework developed here provides a foundation for understanding how neural circuits can maintain function despite structural changes, with potential applications to other neural systems exhibiting homeostatic plasticity.

\appendices



\end{document}



To characterize the space of output-preserving parameter changes, 
we compute the Hessian matrix $H$ at an optimized parameter set $\theta^*$ using automatic differentiation. 
Due to the use of Anderson Acceleration in the fixed-point solver, which employs custom 
reverse-mode differentiation ($custom_vjp$), we cannot use forward-mode autodiff. 
Instead, we compute the Hessian using reverse-mode Jacobian: 
$H = \nabla_\theta (\nabla_\theta y(\theta^*)) = \text{jacrev}(\text{grad}(y))(\theta^*)$, 
where $y(\theta)$ is the output sum function. This approach avoids the forward-mode autodiff 
limitation while maintaining computational efficiency through JAX's JIT compilation.

We then perform an eigendecomposition: $H = U \Lambda U^T$, where $\Lambda$ is a diagonal matrix of eigenvalues $\lambda_i$ and $U$ contains the corresponding eigenvectors $u_i$. The eigenvalue spectrum of the Hessian, computed for the first optimized parameter set, reveals a striking pattern: as shown in Figure~\ref{fig:e_values_of_H}, there are multiple eigenvalues that are extremely small (near zero), with approximately 13 out of 30 eigenvalues having magnitude less than $10^{-6}$. This is highly promising, as it indicates the existence of a substantial manifold of parameter changes along which the output remains approximately constant.

The eigenvectors corresponding to these near-zero eigenvalues define the directions in parameter space where compensatory changes can occur. To understand the structure of these compensatory mechanisms, we examine the eigenvectors associated with the two smallest eigenvalues (closest to zero), shown in Figure~\ref{fig:last_2_e_vector}. These eigenvectors, which represent $\Delta \theta$ directions that preserve output, reveal coordinated patterns of parameter changes. For instance, positive components in certain synaptic weights (indicating increases) are systematically paired with negative components in other weights (indicating decreases), suggesting that increased activity in one pathway must be balanced by decreased activity in another to maintain constant output.

\begin{figure}[H]
    \centering
    \includegraphics[width=0.45\textwidth]{Figures/e_values_of_H.png}
    \caption{Eigenvalue spectrum of the Hessian matrix $H$ computed at the first optimized parameter set. The semilogarithmic plot reveals multiple eigenvalues near zero (below the $10^{-6}$ threshold), indicating directions in parameter space where changes preserve output. Approximately 13 out of 30 eigenvalues are near-zero, suggesting substantial redundancy in the parameter space.}
    \label{fig:e_values_of_H}
\end{figure}

\begin{figure}[H]
    \centering
    \includegraphics[width=0.45\textwidth]{Figures/last_2_e_vectors.png}
    \caption{Eigenvectors corresponding to the two smallest eigenvalues of the Hessian matrix. These eigenvectors represent $\Delta \theta$ directions that preserve output. The pattern of positive and negative components reveals compensatory relationships: increases in certain parameters (e.g., specific synaptic weights) are balanced by decreases in others, maintaining overall circuit function.}
    \label{fig:last_2_e_vector}
\end{figure}

To quantify the effective dimensionality of the output-relevant parameter space across multiple solutions, we computed the participation ratio $PR = (\sum_i \lambda_i)^2 / \sum_i \lambda_i^2$ for 200 optimized parameter sets. The participation ratio measures how many eigenvalues contribute significantly to the output sensitivity: a low participation ratio indicates that only a few eigenvalue directions strongly affect the output, leaving many directions available for compensatory changes. The histogram of participation ratios across all samples, shown in Figure~\ref{fig:ppr}, reveals an unexpected distribution: the majority of samples exhibit low participation ratios, with the histogram peak shifted to the left (low values). This finding is counterintuitive from an optimization perspective, where one might expect most solutions to have high participation ratios (many eigenvalues contributing to output changes). Instead, the left-skewed distribution indicates that most optimized parameter sets lie in regions of parameter space with many near-zero eigenvalues, meaning that most solutions have numerous directions available for output-preserving parameter changes.

\begin{figure}[H]
    \centering
    \includegraphics[width=0.45\textwidth]{Figures/ppr_histogram.png}
    \caption{Histogram of participation ratios computed for 200 optimized parameter sets. The left-skewed distribution (peak at low values) indicates that most solutions have low effective dimensionality, meaning many eigenvalues are near-zero. This suggests that most optimized parameter sets lie in regions with substantial redundancy, providing many directions for compensatory parameter changes that preserve output.}
    \label{fig:ppr}
\end{figure}

This observation has important biological implications. 
The prevalence of low participation ratios suggests that the optimization process naturally 
finds solutions in regions of parameter space with high redundancy. This redundancy is 
precisely what enables homeostatic plasticity: when structural changes occur 
(e.g., during sensory deprivation), the system has multiple compensatory mechanisms 
available to maintain function. The fact that most solutions exhibit this property suggests 
that biological circuits may be naturally configured to operate in such redundant regions, 
facilitating adaptive responses to perturbations.

The next challenge is to analyze the specific structure of these compensatory 
changes and interpret them in biological terms. By examining which parameters i
ncrease and which decrease along the eigenvectors corresponding to near-zero eigenvalues, 
we can identify specific compensatory mechanisms. For example, if we observe increased weights 
in OSN-to-red neuron connections paired with decreased weights in OSN-to-green connections, 
this suggests a compensatory mechanism where reduced input to one pathway is balanced by increased 
input to another. Such patterns could be validated against experimental observations during sensory 
deprivation, where specific structural changes are known to occur.

To better understand these compensatory relationships, future analysis should: 
(1) map the eigenvector components back to specific biological parameters 
(e.g., synaptic weights between specific neuron types), 
(2) identify the most common compensatory patterns across multiple solutions, 
(3) relate these patterns to known mechanisms of homeostatic plasticity such as synaptic 
scaling or structural plasticity, and 
(4) validate predictions by comparing with experimental measurements of parameter changes 
during sensory deprivation. Additionally, a more detailed visualization showing the relationship 
between parameter groups (e.g., all OSN weights, all inter-neuron weights) in the eigenvectors 
would help identify which types of connections are most involved in compensation.

\section{Conclusion}

This technical milestone report presents a mathematical framework for understanding how neural circuits 
can maintain functional output despite structural changes, addressing the fundamental question of 
homeostatic plasticity in the olfactory bulb. We developed a simplified but biologically 
motivated model of the OB circuit, incorporating key neuron types and their connectivity patterns, 
and used the Naka-Rushton activation function to capture nonlinear neural responses.

Through gradient-based optimization using the Adam algorithm, we successfully identified 
multiple parameter configurations that produce target firing rates, 
demonstrating the existence of parameter redundancy—a necessary condition for homeostatic plasticity. Stability analysis confirmed that many of these solutions are dynamically stable, with all eigenvalues of the stability matrix having negative real parts, making them biologically plausible operating points.

The key finding of this work is that the Hessian matrix of the output function exhibits a s
ubstantial number of near-zero eigenvalues across most optimized solutions. 
This indicates the existence of a manifold of parameter changes that preserve output, 
providing a mathematical foundation for understanding compensatory mechanisms. The participation 
ratio analysis revealed that most optimized parameter sets lie in regions of high redundancy, 
with low effective dimensionality of the output-relevant parameter space. This suggests that
biological circuits may naturally operate in such redundant regions, facilitating adaptive responses 
to perturbations.

The eigenvectors corresponding to near-zero eigenvalues encode the 
structure of compensatory changes: coordinated increases and decreases in different parameters that 
maintain overall circuit function. These patterns provide testable predictions about 
which specific parameters should change during sensory deprivation to maintain output, 
offering a bridge between mathematical analysis and experimental validation.

\subsection{Limitations and Future Work}

Several limitations should be acknowledged. First, the model is highly simplified, 
grouping neuron types and neglecting temporal dynamics and oscillatory behavior that are 
characteristic of the real olfactory bulb. Second, we do not include granule cells in the main analysis, 
which play important roles in circuit dynamics. Third, the steady-state assumption ignores important temporal 
aspects of neural responses. Fourth, the biological interpretation of the Hessian eigenvectors 
requires further validation against experimental data on parameter changes during sensory deprivation.

Future work should: (1) incorporate more detailed biological constraints 
(e.g., excitatory-inhibitory balance ratios, known connectivity patterns), 
(2) extend the model to include temporal dynamics and oscillations, (3) v
alidate predicted parameter changes against experimental measurements during sensory deprivation,
 (4) develop methods to map eigenvector patterns to specific biological mechanisms (e.g., synaptic scaling, 
 structural plasticity), and (5) investigate how these mechanisms relate to known forms of homeostatic plasticity. Additionally, analyzing the relationship between participation ratio and other solution properties (e.g., stability, convergence speed) could provide further insights into why optimization finds solutions in redundant regions.

The framework developed here provides a foundation for understanding how neural circuits can maintain 
function despite structural changes, with potential applications to other neural systems exhibiting 
homeostatic plasticity. The mathematical tools—Hessian analysis, eigendecomposition, and participation 
ratio—offer a systematic approach to identifying compensatory mechanisms that could be applied to more
complex models or experimental data.
